%% Chapter 17: Kernelization
%% Status: skeleton -- learning goals and section outline fixed, prose to be written.

\chapter{Kernelization}
\label{ch:kernelization}

Kernelization is preprocessing with a guarantee. We already saw a quadratic
kernel for \lang{Vertex Cover}; here we push it to linear using linear
programming, develop the general tools, and then show that for some problems no
polynomial kernel exists at all.

\begin{goals}
  \poolgoal{19.1}{define \lang{Integer Linear Programming} and show how \lang{Vertex Cover} can be encoded as one}
  \goalitem{prove that the LP relaxation of \lang{Vertex Cover} is half-integral}
  \poolgoal{19.2}{show that \lang{Vertex Cover} has a linear sized kernel}
  \goalitem{use crown decompositions and the sunflower lemma to obtain kernels}
  \goalitem{explain how composition arguments rule out polynomial kernels}
\end{goals}

\section{Integer linear programming}
\label{sec:ilp}

\todo{Write this section.}
% Planned content: The problem, and \lang{Vertex Cover} as an ILP.

\section{Half-integrality and Nemhauser--Trotter}
\label{sec:nemhauser-trotter}

\todo{Write this section.}

\section{A linear kernel for vertex cover}
\label{sec:linear-vc-kernel}

\todo{Write this section.}

\section{Crown decompositions}
\label{sec:crown-decompositions}

\todo{Write this section.}

\section{The sunflower lemma}
\label{sec:sunflower}

\todo{Write this section.}

\section{Lower bounds for kernels}
\label{sec:kernel-lower-bounds}

\todo{Write this section.}
% Planned content: Composition, OR-distillation, and what they rule out.

\section{Exercises}
\label{sec:ex-kernelization}

\todo{Write this section.}
